\chapter{Application Overview}
\label{chap:ch2}

\par To make the comparison easier to follow, since the whole thesis is based on the two backends' code, it will be helpful to see first the Gentlix Bank app as a complete system. This chapter describes the application from the ground up: the banking scenario and features, the React frontend, the database schema, and the target deployment on AWS. A key design goal is that the user always gets the same experience, regardless of which backend is running behind the scenes. That transparency is what makes the comparison meaningful.

\section{Banking scenario and goals}
\label{sec:ch2sec1}

\par Gentlix Bank simulates a simple online banking system. The main goal is 
to let a user authenticate, view account information, and perform basic 
operations on accounts and transactions. The transaction layer is particularly 
data-heavy, since it involves a large volume of records being processed and 
stored, and this makes it a strong point of interest for the comparison. The 
application does not try to cover all real-world banking rules, but it enforces 
enough constraints to be realistic for testing backend behaviour and measuring 
efficiency under load.

\subsection{Functional features}
\label{sec:ch2sec1sub1}

\par From a functional point of view, the system offers the following features:
\begin{itemize}
  \item \textbf{Authentication} --- secure login and registration using email and 
  password, with Argon2id hashing and JWT sessions. One of the 
  main challenges is cross-backend token compatibility: a token issued by the 
  Rust backend must be accepted by the C backend for the same user, and the 
  other way around, because both share one database and are intended to run 
  behind the same load balancer in production.
  \item \textbf{Account management} --- each user has a main account and can 
  create additional banking accounts (savings, credit, and similar types). 
  The domain enforces at most one account per \texttt{(account\_type, currency)} 
  pair, so the user cannot accidentally open duplicate wallets in the same currency.
  \item \textbf{Affiliates} --- users can save other bank members as affiliates 
  with nicknames and quick transfer targets. An affiliate acts as a shortcut: 
  validated account details are already stored, which speeds up sends to frequent 
  recipients.
  \item \textbf{Transactions} --- deposits, withdrawals, sends to another user's 
  account, transfers between the user's own accounts, and card-style payments 
  with merchant metadata. Money is stored as integer cents end-to-end; floating 
  point is not used in the ledger. Every money-moving operation runs inside a 
  single SQL transaction so balance updates and the inserted row commit or roll 
  back together.
  \item \textbf{Dashboard} --- after login, the user sees recent activity and a 
  simple spending chart. Both backends expose a dashboard service that aggregates 
  data server-side so the React app stays thin.
  \item \textbf{Statements and analytics} --- the statement view replays transaction 
  history per account to compute a running \texttt{balance\_after\_cents} on each 
  line, with server-side date filtering, type filtering, sorting, and pagination. 
  Monthly analytics aggregate thousands of rows per user with parallel workers in 
  both backends (see Chapter~\ref{chap:ch5}).
  \item \textbf{REST API} --- all product behaviour is exposed through 22 authenticated 
  routes under \texttt{/api}, plus an unauthenticated \texttt{GET /health} liveness 
  probe. Chapter~\ref{chap:ch6} lists the contract in detail.
\end{itemize}

\subsection{Non-functional requirements}
\label{sec:ch2sec1sub2}

\par Beyond the functional features listed above, the system must meet several 
non-functional requirements that are equally important for the comparison:

\begin{itemize}
    \item \textbf{Security} --- passwords must be stored as secure hashes using 
    a modern algorithm (Argon2id with identical parameters in both backends). 
    JWT tokens must be signed with a shared secret and validated consistently.
    \item \textbf{Consistency} --- both backends must return identical HTTP 
    responses for the same input, including the same status codes, error 
    messages, and JSON field names.
    \item \textbf{Portability} --- the entire application must run in a 
    containerised environment using Docker, so that it can be deployed 
    identically on a developer machine or on AWS.
    \item \textbf{Interoperability} --- a JWT token issued by one backend must 
    be accepted and validated by the other when both run behind the same 
    load balancer and share the same user database.
\end{itemize}

\par The technical stack is as follows: React frontend, C and Rust backends on 
port~6767, PostgreSQL in Docker, JSON over HTTP. Docker Compose in the repository 
supports local development; \texttt{deploy/AWS.md} documents the target cloud layout.

\section{User interface}
\label{sec:ch2sec2}

\par The React frontend focuses not on a large number of pages, but on a clean and practical organisation. The design aims to follow modern standards both visually and structurally, with consistent buttons, forms, and navigation elements throughout the application. Error and success messages are displayed in a clear and immediate way, so that the user always knows the result of their action. Small animations and visual effects have also been added to improve the overall experience.

\par In terms of application flow, the entry point is a landing page with two options: login to access an existing account, or register to open a new one. Each option leads to a dedicated form with the relevant fields. Once authenticated, the user reaches the dashboard, which provides a general summary of their account, along with navigation buttons to the dedicated pages for accounts, affiliates, and transactions. The dashboard also includes a list of recent transactions and a simple chart showing recent spending.

\par The following subsections provide a brief visual overview of the three most important pages in the application.

\subsection{Authentication screen}
\label{sec:ch2sec2subsec1}

\par The authentication screen is the entry point into Gentlix Bank. It contains a simple form with email and password fields and a button to start the login process. If authentication fails, the frontend displays an error message without revealing internal details about the backend or the database.

\begin{figure}[H]
    \centering
    \includegraphics[width=0.9\textwidth]
    {figures/auth_screen.jpeg}
    \label{fig:auth_screen}
\end{figure}

\subsection{Accounts dashboard}
\label{sec:ch2sec2subsec2}

\par The dashboard is designed to give the user a quick overview without overwhelming them with information. Each account information is displayed as a separate card. Clicking on a card opens the detailed page for that feature. The recent transactions list and the spending chart at the bottom of the page update automatically to reflect the user's activity.

\begin{figure}[H]
    \centering
    \includegraphics[width=0.9\textwidth]
    {figures/dashboard_screen.jpeg}
    \label{fig:dashboard_screen}
\end{figure}

\subsection{Transaction details and creation}
\label{sec:ch2sec2subsec3}

\par The transaction view shows recent operations for a selected account. It also offers forms to create deposits, withdrawals, transfers, sends, and payments. The frontend validates simple rules (for example, non-empty fields and positive amounts) before sending the request to the backend. Heavy work---running balances, filters, and analytics---stays on the server so both backends can be compared fairly under load.

\section{Database schema and domain model}
\label{sec:ch2sec3}

\par At the core of Gentlix Bank there is a shared PostgreSQL database used by 
both backends. The database schema is intentionally simple, but it still 
captures the main entities needed for an online banking scenario: users, 
accounts, affiliates, and transactions. Each entity maps directly to one of 
the functional features described earlier in this chapter.

\par The \texttt{users} table stores the basic identity information and login 
credentials for each customer. The \texttt{accounts} table represents the 
different banking accounts owned by a user, such as the main account, savings 
accounts, or credit accounts. The \texttt{affiliates} table stores shortcuts 
to other users in the system, so that transfers to frequent recipients can be 
created faster. Finally, the \texttt{transactions} table records all money 
movements between accounts.

\begin{figure}[htbp]
    \centering
    \includegraphics[width=0.8\textwidth]
    {figures/db_schema.png}
    \label{fig:db_schema}
\end{figure}

\par Relationships between tables reflect the real-world constraints of the 
banking model. Each account belongs to exactly one user, and each transaction 
references both a source and a destination account. Affiliates link a user to 
another user and store the information needed to start a transfer without 
re-entering all details. This structure ensures that both the C and Rust 
backends work with the same domain model and can be compared fairly.

\par The database schema is shared by both implementations through migration 
files \texttt{001}--\texttt{003} in the repository. A benchmark seed script 
loads roughly 100\,000 transactions for one test user so that statement and 
analytics workloads in Chapter~\ref{chap:ch8} stress realistic data volumes.

\section{Target deployment on AWS}
\label{sec:ch2sec4}

\par Gentlix Bank is designed for a small but representative cloud layout on 
Amazon Web Services (AWS). In production, the React frontend is served as static 
files, and HTTP API traffic goes through a load balancer to two EC2 instances: 
one hosting the C backend and one hosting the Rust backend. Both instances 
connect to the same Dockerised PostgreSQL database and receive the same 
\texttt{DATABASE\_URL} and \texttt{JWT\_SECRET} through environment variables.

\par The two instances should use the same instance type and similar operating 
system configuration so that any future cloud benchmarks are not skewed by 
hardware differences. From the user's perspective, the experience stays 
transparent: the browser talks to one URL and does not know which language 
handled a given request.

\par \textbf{Local development and benchmarks.} Day-to-day development and the 
k6 measurements reported in Chapter~\ref{chap:ch8} use \texttt{docker-compose} 
or a native backend on \texttt{localhost:6767} with shared PostgreSQL. That setup 
keeps network noise out of the comparison while preserving the same schema and 
API. AWS remains the documented deployment target; repeating the benchmarks in 
the cloud is optional follow-up work if time allows.
